\chapter{Discussion}
\label{chap:discussion}

This thesis asked whether the Welch--Satterthwaite principle can be adapted
to compare mutual information between two independent categorical
populations. The central problem is not estimation of the MI difference
itself. It is calibration of the standardized difference when the two
first-order variance contributions are estimated from finite, heterogeneous
contingency tables. This chapter interprets the proposed solution and the
evidence obtained for it.

\section{Answers to the research questions}

\subsection{Method construction}

Research Question~1 asked how to adapt Welch--Satterthwaite inference to the
nonlinear plug-in MI variance estimator. The expanded method does this in two
stages. For each population, it first measures the sampling variability of
the estimated MI variance through the influence function $g_P$ or $g_Q$.
It then matches the mean and variance of that positive variance estimator to
a scaled chi-squared variable. This produces the component degrees of freedom

\begin{equation}
  \widehat\nu_P
  =\frac{2n_P\widehat V(P)^2}{\widehat\tau^2(P)},
  \qquad
  \widehat\nu_Q
  =\frac{2n_Q\widehat V(Q)^2}{\widehat\tau^2(Q)}.
  \label{eq:discussion-component-df}
\end{equation}

The ordinary Welch--Satterthwaite equation combines these values according to
the two contributions to the squared standard error. The result is a Student
reference whose tail weight is determined by the observed stability of the
MI variance estimates. No resampling or numerical optimisation is required at
test time.

The construction is MI-specific because $g_P$ differentiates the complete
variance functional. In particular, it includes the changes in pointwise MI
caused by changing the joint and marginal probabilities. Replacing this step
by $n_P-1$, as in a direct analogy with Welch's test for means, does not
represent the sampling uncertainty of the MI variance estimator.

\subsection{Statistical performance}

Research Question~2 asked whether expanded Welch improves type-I error and
coverage. Across all 60 equal-MI population pairs, it reduced mean absolute
false-positive-rate error at $\alpha=0.05$ from 0.01342 for normal Wald and
0.01194 for simple Welch to 0.00825. Overall 95\% coverage increased from
0.93799 for normal Wald to 0.94455.

The improvement was concentrated in the regimes for which the method was
designed. Expanded Welch reduced calibration error by 29.3--46.5\% relative
to normal Wald in the highly skewed, ultra-skewed, and widespread-sparsity
regimes. It also improved the moderate regime by 37.8\%. In the well-sampled
control, normal Wald was already accurate and expanded Welch was effectively
similar, although its mean absolute error was 4.5\% larger. The evidence
therefore supports improved average calibration in difficult positive-MI
settings, not uniform dominance for every population pair.

The most extreme scenario remained outside satisfactory calibration. Under
ultra-sparse cells and a $10:1$ sample-size ratio, expanded Welch reduced
mean FPR from 0.09233 to 0.07343 but did not reach the nominal value of 0.05.
The reference correction addresses uncertainty in the estimated variance; it
does not remove all finite-sample bias, discreteness, or empirical support
loss.

\subsection{Practical trade-offs}

Research Question~3 asked about power, runtime, validity, and operating
range. The heavier Student tails reduced power by 1.02 percentage points on
average across the five alternatives and by at most 1.23 points. This is a
small but genuine cost. A liberal null reference can appear more powerful
partly because it rejects too often under both the null and alternative. The
calibration and power results must therefore be interpreted together.

Expanded Welch took approximately 0.15 milliseconds per table pair in the
timed implementation, compared with 0.08 milliseconds for normal Wald. Both
methods have $O(rc)$ time and memory complexity. The correction is therefore
about 1.7 times slower in the tested implementation but remains negligible in
absolute cost and avoids the multiplicative cost of resampling.

The valid rate was effectively one outside widespread sparsity. It fell to
0.97447 when support loss caused the empirical first-order variance to vanish
in some replicates. A finite p-value is not the only criterion for a useful
test; the underlying first-order variance must also be identifiable from the
observed table.

\section{Why the correction improves difficult cases}

Normal Wald treats the estimated standard error as if it were known. This is
a reasonable approximation when the table is well sampled, because repeated
samples then produce similar empirical probabilities and similar variance
estimates. Skewness, rare cells, and unequal sample sizes make the denominator
of the test statistic less stable. A normal reference does not respond to
that additional uncertainty.

Simple Welch adds finite degrees of freedom, but assigns them from sample
sizes alone. Its reference changes little when two tables have the same
sample sizes but very different cell probabilities. The primary results show
the consequence: simple Welch generally lies between normal Wald and expanded
Welch, but its improvement is modest in the most difficult regimes.

Expanded Welch estimates how strongly the complete MI variance functional
changes when probability is moved towards each cell. Large variation among
the resulting $g_P(i,j)$ values implies a volatile variance estimate, a
smaller component degree of freedom, and a heavier-tailed Student reference.
The correction therefore becomes strongest when the observed table indicates
that the standard error itself is unreliable. As sampling improves,
$\widehat\nu_P$ and $\widehat\nu_Q$ increase and the Student reference
approaches the normal reference.

This mechanism also explains the residual failures. The derivation is
first-order. When many cells disappear from the empirical support, both the
point estimate and its estimated variance can be governed by higher-order and
discrete effects. Changing only the reference tails cannot reconstruct
information that is absent from the table.

\section{Interpretation of the scaled chi-squared model}

Satterthwaite's scaled chi-squared representation is a working model for the
positive estimated variance, not an assertion that its exact finite-sample
distribution is chi-squared. Its role is to translate the first two moments
of the variance estimator into an effective degree of freedom. That degree
of freedom then enters the final Student reference.

The separate mechanism audit supports this use with qualifications. An
oracle scaled chi-squared distribution fitted the complete distribution about
as well as an oracle normal model and better than an oracle lognormal model.
It produced the smallest upper-tail error among the three candidates. The
larger discrepancy arose in the first-order prediction of finite-sample
moments, not in the choice between chi-squared and normal shape families.
The plug-in component degrees of freedom nevertheless had median ratio 0.970
to empirical moment-matched values.

These findings justify scaled chi-squared moment matching as a useful
approximation for tail calibration. They do not establish an exact pivot.
The finite-sample performance demonstrated in Chapter~\ref{chap:results}
remains empirical evidence for a theoretically motivated approximation.

\section{Relation to existing methods}

The expanded procedure occupies a middle position between normal Wald
inference and resampling. Like Wald inference, it is deterministic, analytic,
and linear in the number of cells. Like Welch--Satterthwaite inference, it
acknowledges that an estimated denominator changes the reference tails. Its
distinctive step is to calculate that denominator uncertainty for the MI
variance functional rather than importing the ordinary sample-variance
degrees of freedom.

The method also differs from an independence test. Pearson and likelihood-
ratio tests compare a single joint distribution with the product of its own
margins and use a multivariate quadratic statistic. The present method
compares one scalar functional across two independent populations and uses a
signed, studentized difference. At the independence boundary, the first-order
MI influence vanishes and the regular derivation used here no longer applies.
Appendix~\ref{app:independence-boundary} develops this distinction.

The literature review identified prior ingredients rather than a direct
precedent for the complete method. Bias correction for information measures,
first-order MI variance, influence-function analysis, tests for entropy-based
indices, and Welch--Satterthwaite moment matching are established. The claim
of this thesis is consequently methodological and specific: it constructs and
evaluates an MI-variance influence-function degree-of-freedom correction for
the two-sample equal-MI problem. It does not claim invention of the component
theories.

\section{Limitations}

Several limitations define the strength of the conclusions.

First, the primary evidence is simulation-based. The 60 fixed population
pairs cover multiple table sizes, marginal shapes, association structures,
sample ratios, and expected-count regimes, but they do not span all discrete
distributions. A new independently generated confirmatory grid would provide
a stronger guard against conclusions that depend on the chosen population
generator.

Second, the theory assumes fixed finite alphabets, positive population
support, independent samples, and positive first-order MI variance. The
primary experiment deliberately uses $I(P)=I(Q)>0$. It does not validate
the method at independence, where the first derivative is zero, or when the
alphabet grows with sample size.

Third, every analytic method uses the same first-order standard error and the
same Miller-type bias correction. Expanded Welch changes the reference
distribution only. Residual errors can therefore arise from finite-sample
bias in the numerator, error in the first-order variance estimate, dependence
between numerator and denominator, or discreteness of the table distribution.
The experiment does not identify one of these as the sole remaining cause.

Fourth, invalid replicates are conditioned out uniformly across methods. This
is appropriate for comparing references on common inputs, but a low valid
rate would limit practical use regardless of conditional calibration. The
widespread-sparsity result illustrates this point.

Fifth, the power experiment contains five controlled alternatives rather than
a broad power surface. It establishes the approximate cost of the heavier
reference in representative $3\times3$ cases, but does not characterise
power over all effect sizes, table shapes, or sample ratios.

Finally, no real-data case study is included. The thesis establishes a method
and its controlled numerical behaviour. Applied work is needed to determine
how often its difficult regimes arise in specific fields and whether its
calibration advantage changes scientific conclusions.

\section{Recommended operating range}

The results support expanded Welch when two independent categorical samples
have positive estimated MI and the analyst seeks a fast analytic comparison
of their population MI values. It is most useful when one or both tables have
skewed margins, rare expected cells, or unequal sample sizes, but retain enough
observed support for positive first-order variance estimates.

Normal Wald remains adequate in well-sampled tables and is simpler. Expanded
Welch should not be presented as a repair for arbitrary sparsity. If complete
rows or columns are frequently absent, if the estimated first-order variance
is zero, or if the scientific null is independence, the regular two-sample
method is not the appropriate reference without further development.

\section{Further work}

The first priority is an independently generated confirmatory experiment with
the analysis code and decision rules fixed in advance. It should preserve the
current regime definitions, report scenario-level uncertainty, and include a
broader power surface. This would test whether the observed calibration gains
generalise beyond the development grid.

The second priority is improved finite-sample centring. The expanded method
models uncertainty in the denominator but does not account for all
distribution-specific bias in the difference. An analytically tractable
higher-order correction could address the remaining liberal behaviour without
introducing resampling.

The third priority is boundary theory. Near $I=0$, first-order normal
inference transitions to the second-order chi-squared geometry of
independence testing. A unified method would require an explicit treatment of
that transition rather than a direct extension of the present degrees of
freedom.

Finally, the same variance-functional strategy may be investigated for
conditional MI and related discrete information measures. Such extensions
should follow, rather than precede, independent validation of the simpler
two-sample MI setting.
